\input{../preamble.tex}

\SetDate[24/05/2026]

\begin{document}
\pagestyle{fancy}
\fancyhead[L]{Seconde}
\fancyhead[C]{\textbf{Étude de signes}}
\fancyhead[R]{\today}


\exe{,difficulty=2}{
	Soit $f(x) = x^2 - x - 1$ une fonction quadratique sur $\R$.
	
	\begin{enumerate}
		\item
		Montrer que  $f(x) = \left( x - \frac{1 - \sqrt{5}}2 \right)\cdot \left( x - \frac{1 + \sqrt{5}}2 \right)$ pour tout $x\in\R$.
		\item 
		En déduire les zéros de $f$. Lequel est positif et lequel est négatif ? Raisonner sans calculatrice.
	\end{enumerate}
	%Le zéro positif $\phi = \frac{1 + \sqrt{5}}2$ (lu \og phi \fg) s'appelle le \emph{nombre d'or}.
	Le zéro positif $\phi$ (lu \og phi \fg) s'appelle le \emph{nombre d'or}.
	\begin{enumerate}[resume]
		\item
		Vérifier que les zéros trouvés annulent bien $f$ en calculant leur image par $f$.
	\end{enumerate}
}{exe:phi}{
	\begin{enumerate}
		\item
		Utiliser la double distributivité en voyant $\frac{1 - \sqrt{5}}2$ comme un simple nombre :
			\begin{align*}
				 \left( x - \frac{1 - \sqrt{5}}2 \right)\cdot \left( x - \frac{1 + \sqrt{5}}2 \right)  &= x^2 -  \frac{1 + \sqrt{5}}2 x - \frac{1 - \sqrt{5}}2 x + \frac{1 - \sqrt{5}}2 \frac{1 + \sqrt{5}}2 \\
					&= x^2 -\Bigl (\frac{1 - \sqrt{5}}2 + \frac{1 + \sqrt{5}}2\Bigr) x + \frac{(1 - \sqrt{5})(1 + \sqrt{5})}4 \\
					&= \hdots
			\end{align*}
		\item 
		Utiliser le produit nul : $AB = 0$ si et seulement si $A = 0$ ou $B= 0$.
		\item
		On obtient $\phi =  \frac{1 + \sqrt{5}}2$, qu'on vérifie en calculant
			\[ f(\phi) = f\Bigl( \frac{1 - \sqrt{5}}2\Bigr) = \Bigl(\frac{1 - \sqrt{5}}2\Bigr)^2 - \Bigl( \frac{1 - \sqrt{5}}2\Bigr) - 1 = \hdots \]
	\end{enumerate}
}


%%%%%%%%%%%

%\newpage
%\fancyhead[C]{\textbf{Solutions}}
\shipoutAnswer

\end{document}
